% !TEX root = ../supporting_information.tex

\SISection{Field-scale reservoir simulation of geologic
\texorpdfstring{\coTwo{}}{CO2} storage}
\label{sec:si-field-scale-reservoir-simulation}

\SISubsection{Reservoir model and simulation setup}

The field-scale model contains approximately two million active cells and
represents megatonne-scale \coTwo{} injection into the offshore Texas storage
interval followed by migration and containment over 1,000 years. JutulDarcy
solves the multiphase flow problem. Each realization preserves the common
reservoir geometry and operating conditions while replacing the geology-
specific stratigraphy and the spatially variable fault permeability,
porosity, capillary-pressure functions, and relative-permeability functions.

\todo{Document the JutulDarcy version, grid and fluid model, injection
schedule, well controls, and boundary and initial conditions.}

\SISubsection{Reservoir-input validation}

Before model construction, the importer verifies that the fault-property and
stratigraphy files share the same geology identifier, design metadata, and
compatible hashes. It then checks complete six-window by 87-slice coverage,
permeability units and coordinate convention, porosity bounds, capillary-
pressure and relative-permeability region indices, and consistency of the
native saturation endpoints.

\SISubsection{Numerical verification and restart procedures}

A case is accepted for analysis only after the simulation reaches the
requested final time and passes the prescribed mass-balance and output-
integrity checks.

\todo{Document the nonlinear and linear solvers, timestep controls, restart
procedure, mass-balance tolerance, five saved output times, failed-case
recovery rules, and final validation results.}
